% \newtheorem{introtheorem}{Theorem}
% \renewcommand{\theintrotheorem}{\Alph{introtheorem}} % "letter-numbered" theorems

% \swapnumbers
% \theoremstyle{plain}
% \newtheorem{theorem}{Theorem}[section]
% \newtheorem{lemma}[theorem]{Lemma}
% \newtheorem{claim}[theorem]{Claim}
% \newtheorem{corollary}[theorem]{Corollary}

% \theoremstyle{remark}
% \newtheorem{remark}[theorem]{Remark}
% \newtheorem*{remark*}{Remark}
% \newtheorem{example}[theorem]{Example}

% \theoremstyle{definition}
% \newtheorem{definition}[theorem]{Definition}
% \newtheorem{miniremark}[theorem]{}

% Hom(R^n,R^n)
% \newcommand{\End}[1]{ \mathcal{Z}_{#1} }
% \newcommand{\Endd}[1]{ \mathcal{Z}^*_{#1} }
\newcommand{\End}[1]{ \mathrm{End}({#1}) }
% \newcommand{\EndR}[1]{ \End{\R^{#1}} }
\newcommand{\Endd}[1]{ \mathrm{End}({#1})^* }
% \newcommand{\EndRd}[1]{ \Endd{\R^{#1}} }
% \newcommand{\End}[1]{ \mathrm{End}(#1) }
% \newcommand{\Endd}[1]{ \mathrm{End}(#1)^* }

\newcommand{\SymEnd}[1]{ \mathrm{Sym}({#1}) }
\newcommand{\AsymEnd}[1]{ \mathrm{Asym}({#1}) }

%
\newcommand{\GF}[1]{ \mathcal{G}_{#1} }
\newcommand{\GFd}[1]{ \mathcal{G}^{*}_{#1} }

% Sobolev spaces of varifolds
\newcommand{\VSobnorm}[3]{\mathbf{W}_{#1} ( {#2}, {#3} )}
\newcommand{\VSobloc}[1]{\mathbf{W}_{#1}^{\mathrm{loc}}}
\newcommand{\VSob}[1]{\mathbf{W}_{#1}}
\newcommand{\VSobz}[1]{\mathbf{W}^{\diamond}_{#1}}

% the L^q norm of a function u in U is '\norm{u}{q}{A}' for "|u|_{q;U}" with
% bold |
\newcommand{\norm}[3]{\boldsymbol{|} #1 \boldsymbol{|}_{#2;#3}}

% \Shuffle(p,q)
% Shuffle of type (p,q), i.e., the set of permutations of {1,2,...,p+q}
% which are increasing on each of the sets {1,...,p} and {p+1,...p+q};
% see 1.3.3 in Federer1969.
\newcommand{\Shuffle}[2]{\operatorname{\textbf{Sh}}(#1,#2)}

% \Choice(n,k)
% All choices of k elements from the set of n elements realised
% as the set of all increasing maps from {1,...,k} to {1,...,n};
% see 1.3.2 in Federer1969.
\newcommand{\Choice}[2]{\boldsymbol{\Lambda}(#1,#2)}

% \Partition(k,n) (a.k.a. multi-indices on \R^k of length n)
% Partitions of n into k subsets realised as the set of functions
% \alpha mapping the set {1,...,k} into non-negative integers
% such that \sum_{i=1}^k \alpha(i) = n ; see 1.9.2 in Federer1969.
\newcommand{\Partition}[2]{\boldsymbol{\Xi}(#1,#2)}

% \Sequence(n,k)
% The set of functions of the type {1,...,k} --> {1,...,n} ;
% see 1.10.6 in Federer1969.
\newcommand{\Sequence}[2]{\mathscr{S}(#1,#2)}

% The space of continuous functions with compact support, "K(X)" with
% calligraphic "K", produced by '\ccspace{X}'
\newcommand{\ccspace}[1]{\mathscr{K}(#1)}

% The space of smooth test functions with compact support
\newcommand{\dspace}[2]{\mathscr{D}(#1,#2)}

% The space of distributions
\newcommand{\distr}[2]{\mathscr{D}'(#1,#2)}

% The space of smooth functions
\newcommand{\espace}[2]{\mathscr{E}(#1,#2)}

% turning a plane into a projection, i.e. obtain '\project{T}' produces
% "T_\natural", the projection associated with the plane T
\newcommand{\project}[1]{{#1}_\natural}

% similar, with parentheses
\newcommand{\eqproject}[1]{(#1)_\natural}

% again similar for the orthogonal complement of the plane T
\newcommand{\perpproject}[1]{{#1}^\perp_\natural}

% the Grassmannian '\grass{n}{m}' produces "G(n,m)" with bold "G", i.e. the
% set of m dimensional planes in R^n
\newcommand{\grass}[2]{\mathbf{G}(#1,#2)}

% The projective Grassmannian,
% i.e., $\pgrass{X}{k}$ is the set of all orthogonal projections in X of rank k
\newcommand{\pgrass}[2]{\mathbf{G}_{\natural}(#1,#2)}

% the oriented Grassmannian
\newcommand{\ograss}[2]{\mathbf{G}_0(#1,#2)}

\newcommand{\AllP}{\mathbf{P}}

% the manifold of all projections
\newcommand{\Proj}[2]{\AllP(#1,#2)}

% The measure of a ball, '\measureball{\mu}{\oball{a}{r}}' produces '\mu
% U(a,r)', i.e. no brackets for measure of balls for readability
\newcommand{\measureball}[2]{{#1}\,{#2}}

% Space of Q valued functions
\newcommand{\qspace}{\mathbf{Q}}

% Commands for varifolds
\newcommand{\Var}[1]{\mathbf{V}_{#1}}     % general varifolds
\newcommand{\RVar}[1]{\mathbf{RV}_{#1}}   % rectifiable varifolds
\newcommand{\IVar}[1]{\mathbf{IV}_{#1}}   % integral varifolds
\DeclareMathOperator{\VarTan}{VarTan}     % varifold tangents
\newcommand{\var}[1]{\mathbf{v}_{#1}}     % varifolds generated by rectifiable sets

\DeclareMathOperator{\graph}{graph}     % the graph of a function

% For ease of changing these letters
\newcommand{\adim}{n}
\newcommand{\vdim}{k}
\newcommand{\codim}{n-k}

% integral currents
\newcommand{\IC}{\mathbf{I}}

% integral currents
\newcommand{\NC}{\mathbf{N}}

% the open ball: '\oball{a}{r}' produces "U(a,r)" with bold "U"
\newcommand{\oball}[2]{\mathbf{U}(#1,#2)}

% the closed ball: '\cball{a}{r}' produces "B(a,r)" with bold "B"
\newcommand{\cball}[2]{\mathbf{B}(#1,#2)}

% the closed ball: '\cball{a}{r}' produces "B(a,r)" with bold "B"
\newcommand{\cballa}[3]{\mathbf{B}^{#1}(#2,#3)}

% the sphere
\newcommand{\sphere}[1]{\mathbb{S}^{#1}}

% the canonical isomorphism between \End{X} and X^* \otimes X
\newcommand{\caniso}[1]{\boldsymbol{\Gamma}_{\!#1} }

% textstyle integral signs
\ifdefined\textint
    \renewcommand{\textint}[2]{{\textstyle\int_{#1}^{#2}}}
\else    
    \newcommand{\textint}[2]{{\textstyle\int_{#1}^{#2}}}
\fi

% textstyle integral signs
\ifdefined\textfint
    \renewcommand{\textfint}[2]{{\textstyle\fint_{#1}^{#2}}}
\else    
    \newcommand{\textfint}[2]{{\textstyle\fint_{#1}^{#2}}}
\fi

% textstyle sum sign
\ifdefined\textsum
    \renewcommand{\textsum}[2]{{\textstyle\sum_{#1}^{#2}}}
\else  
    \newcommand{\textsum}[2]{{\textstyle\sum_{#1}^{#2}}}
\fi

% textstyle product
\ifdefined\textint
    \renewcommand{\textprod}[2]{{\textstyle\prod_{#1}^{#2}}}
\else    
    \newcommand{\textprof}[2]{{\textstyle\prod_{#1}^{#2}}}
\fi


% natural numbers
% \newcommand{\nat}{\mathbb{N}}
% \newcommand{\nat}{\mathbb{N}_+}
\newcommand{\natp}{\mathscr{P}}
\newcommand{\nat}{\natp \cup \{0\}}

% the integers
\newcommand{\integers}{\mathbf{Z}}
\newcommand{\Z}{\integers}

% for 0-dimensional currents: [[ x ]]
% \newcommand{\Lcur}{\boldsymbol{[}}
% \newcommand{\Rcur}{\boldsymbol{]}}
% \newcommand{\Lcur}{[\![}
% \newcommand{\Rcur}{]\!]}
\newcommand{\Lcur}{\llbracket}
\newcommand{\Rcur}{\rrbracket}

% admissible deformations
\newcommand{\adm}[1]{\mathfrak{D}({#1})}

% complex numbers
\newcommand{\C}{\mathbf{C}}

% real numbers
\newcommand{\R}{\mathbf{R}}

% cubical subdivision of R^n
\newcommand{\W}{\mathbf{W}}

% characteristic function
% \newcommand{\CF}{\mathds{1}}
\newcommand{\CF}[1]{\boldsymbol{\chi}_{#1}}

% the Čech homology
\newcommand{\Ch}{\mathbf{\check H}}

% the singular homology
\newcommand{\Sh}{\mathbf{H}}

% homology
\newcommand{\h}{\mathbf{H}}

% cycles for homology theory
\newcommand{\hZ}{\mathscr{Z}}

% boundaries  for homology theory
\newcommand{\hB}{\mathscr{B}}

% the Lebesgue measure
\newcommand{\LM}{\mathscr{L}}

% The Lebesgue space, '\Lp{p}' produces "L_p" with bold L
\newcommand{\Lp}[1]{\mathbf{L}_{#1}}

% The local Lebesgue space, '\Lploc{p}' produces "L_p^{loc}" with bold L
\newcommand{\Lploc}[1]{\mathbf{L}_{#1}^{\mathrm{loc}}}

% Slightly old-fashened notation, may explain or replace
% The L^p-norm of u w.r.t. \phi, '\Lpnorm{\phi}{p}{u}' produces
% "\phi_{(p)}(u)", unusual but efficient. Also
\newcommand{\Lpnorm}[3]{{#1}_{({#2})}({#3})}
\newcommand{\eqLpnorm}[3]{{(#1)}_{({#2})}({#3})}

% the number in Besicovitch's covering theorem in n dimensions,
% '\besicovitch{n}', produces "\beta(n)" with bold beta
\newcommand{\besicovitch}[1]{\boldsymbol{\beta}(#1)}

% The nearest point projection
\newcommand{\npp}[1]{\boldsymbol{\xi}_{{#1}}}

% The distance function
% \newcommand{\dist}[1]{\boldsymbol{\delta}_{{#1}}}
\DeclareMathOperator{\dist}{dist}

% the distance function
\newcommand{\distF}[2]{\boldsymbol{\delta}_{#2}^{#1}}
% \da{A}{\phi} will produce \boldsymbol{\delta}_{A}^{\phi}
\newcommand{\da}[2]{\distF{#2}{#1}}

% nu
\newcommand{\an}[2]{\boldsymbol{\nu}_{#1}^{#2}}

% r
\newcommand{\ar}[2]{\boldsymbol{r}_{#1}^{#2}}

% \xi
\newcommand{\anp}[2]{\boldsymbol{\xi}_{#1}^{#2}}

% boldface p
\newcommand{\pp}{\mathbf{p}}

% boldface q
\newcommand{\qq}{\mathbf{q}}

% The domain of the unique nearest point projection
\newcommand{\Unp}[1]{\mathrm{Unp}({#1})}
% \DeclareMathOperator{\Unp}{Unp}


% % distance
% \DeclareMathOperator{\dist}{dist}

% the Hausdorff measure
\newcommand{\HM}{\mathscr{H}}

% the Hausdorff distance
\newcommand{\HD}{{d_{\HM}}}
% \newcommand{\HD}{\operatorname{HD}}

\newcommand{\HDK}[1]{{d_{\HM,#1}}}

% the flat metric on varifolds
\newcommand{\FM}{\mathbf{L}}

% the density function for measures
\newcommand{\density}{\boldsymbol{\Theta}}
\newcommand{\udensity}{\boldsymbol{\overline{\Theta}}}
\newcommand{\ldensity}{\boldsymbol{\underline{\Theta}}}

\newcommand{\BFconst}{\boldsymbol{\beta}}

\newcommand{\eqclsl}{\boldsymbol{[}}
\newcommand{\eqclsr}{\boldsymbol{]}}

% measure of the n dimensional unit ball, '\unitmeasure{n}' produces
% "\alpha(n)" with bold alpha
\newcommand{\unitmeasure}[1]{\boldsymbol{\alpha}(#1)}

% The restriction of a measure \phi to a set A is '\phi \restrict A'
% \newcommand{\restrict}{\mathop{\llcorner}}
\newcommand{\restrict}{ \mathop{ \rule[0pt]{.5pt}{6pt} \rule[0pt]{4pt}{0.5pt} }\nolimits }
% \newcommand{\restrict}{\lefthalfcup}

% The contraction of a form with a vectorfield
% \newcommand{\contract}{\righthalfcup}
\newcommand{\contract}{ \mathop{ \rule[0pt]{4pt}{0.5pt} \rule[0pt]{.5pt}{6pt} }\nolimits }

% the upright "d" in integrals
\newcommand{\ud}{\ensuremath{\,\mathrm{d}}}

% the upright "D" for differentiation
\newcommand{\uD}{\ensuremath{\mathrm{D}}}
% \DeclareMathOperator{\Der}{D}       
\newcommand{\Der}{\uD}

\DeclareRobustCommand{\rchi}{{\mathpalette\irchi\relax}}
\newcommand{\irchi}[2]{\raisebox{\depth}{$#1\chi$}}

\newcommand{\arl}[2]{\underline{\ar{#1}{#2}}}

\DeclareMathOperator{\cato}{cato}

\DeclareMathOperator{\extreme}{extreme}
\DeclareMathOperator{\exposed}{exposed}


% the weak derivative with bold D instead of D for the classical derivative
% \newcommand{\weakD}{\operatorname{\mathbf{D}}}
\newcommand{\weakD}{\ensuremath{\mathbf{D}}}

% the identity mapping on A, '\id{A}', produces "1_A" with bold 1
% \newcommand{\id}[1]{\mathbf{1}_{#1}}
% \newcommand{\id}[1]{\operatorname{id}_{#1}}
% \newcommand{\id}[1]{\mathrm{id}_{#1}}
% \newcommand{\id}[1]{\mathbf{I}_{#1}}
\newcommand{\id}[1]{\mathds{1}_{#1}}
% \newcommand{\id}[1]{\mathbb{1}_{#1}}

% The image of a set A under a map f is "f[A]"
\newcommand{\lIm}{[}
\newcommand{\rIm}{]}

\newcommand{\scale}[1]{\boldsymbol{\mu}_{#1}}
\newcommand{\trans}[1]{\boldsymbol{\tau}_{#1}}

% the space of a simplicial complex
\newcommand{\sspace}[1]{\textstyle{\bigcup}\,{#1}}

% \newcommand{\tbwedge}{\text{{\small $\bigwedge$}}}
\newcommand{\tbwedge}{{\textstyle\boldsymbol{\bigwedge}}}

% exterior algebra
% \newcommand{\extpower}[1]{\text{\raisebox{2pt}{\scriptsize $\bigwedge$}}_{#1}}
% \newcommand{\extpower}[1]{{\textstyle \bigwedge^{#1}}}

% space of anti-symmetric forms
\newcommand{\aforms}[1]{\tbwedge^{\!{#1}}}
% \newcommand{\aforms}[1]{\bigwedge\nolimits^{\!{#1}}}

% exterior power of a vectorspace
\newcommand{\extpower}[1]{\tbwedge_{{#1}}}
% \newcommand{\extpower}[1]{\bigwedge\nolimits_{{#1}}}

% % exterior algebra
% % \newcommand{\extpower}[1]{\text{\raisebox{2pt}{\scriptsize $\bigwedge$}}_{#1}}
% \newcommand{\extpower}[1]{{\textstyle \bigwedge^{#1}}}
% \newcommand{\tbwedge}{{\textstyle \bigwedge}}

% Adjoint of a linear map
\newcommand{\adjoint}[1]{\aforms{1} {#1}}
\newcommand{\aadjoint}[1]{\aforms{1} \aforms{1} {#1}}

% transpose of a linear map between Euclidean spaces
% \transpose{A} produces A^{*}
\newcommand{\transpose}[1]{ #1^{*} }
\newcommand{\ttranspose}[1]{ #1^{**} }

% convex transpose of a convex map
\newcommand{\ctran}[1]{{#1}^{\dagger}}
\newcommand{\cttran}[1]{{#1}^{\dagger\dagger}}

% polar of a function or a set
\newcommand{\polar}[1]{{#1}^{\circ}}
\newcommand{\ppolar}[1]{{#1}^{\circ\circ}}

% convex conjugate of a extended real valued map
\newcommand{\cconj}[1]{{#1}^{\star}}
\newcommand{\ccconj}[1]{{#1}^{\star\star}}

% textstyle big cup
\newcommand{\tbcup}{{{\textstyle \bigcup}}}

% textstyle big cap
\newcommand{\tbcap}{{{\textstyle \bigcap}}}

% The closure of a set A, '\Clos{A}' produces "Clos A"
\newcommand{\Clos}[1]{\mathop{\mathrm{Clos}}#1}

% The centre, side length, and "lower left" corner of a cube
\newcommand{\side}[1]{\mathbf{l}(#1)}
\newcommand{\centre}[1]{\mathbf{c}(#1)}
\newcommand{\vertex}[1]{\mathbf{o}(#1)}

% The interior and boundary of a cube
\newcommand{\cInt}[1]{\operatorname{Int}_{\mathrm{c}}(#1)}
% \newcommand{\cBdry}[1]{\operatorname{Bdry}_{\mathrm{c}}(#1)}
\newcommand{\cBdry}[1]{\partial_{\mathrm{c}}#1}
% \newcommand{\cInt}[1]{\mathop{\mathbf i}(#1)}
% \newcommand{\cBdry}[1]{\mathop{\mathbf b}(#1)}

% The standard dyadic cubes
\newcommand{\cubes}{\mathbf{K}}

% The Whitney family of cubes
\newcommand{\WF}{\mathbf{WF}}

% The cubical complex of a family of cubes
\newcommand{\CX}{\mathbf{CX}}

% Two commands for cylinders, '\cylind{T}{a}{r}' produces "C(T,a,r)" with bold
% C, for plane T, center a, radius r
\newcommand{\cylind}[3]{\mathbf{C} ( #1, #2, #3 )}

% a variant for cylinders with finite height
\newcommand{\cylinder}[4]{\mathbf{C} ( #1, #2, #3, #4 )}

% A symbol for the Dirac measure
% \newcommand{\Dirac}[1]{\operatorname{Dirac}(#1)}
% ... following the French/Bourbaki school
% \newcommand{\Dirac}[1]{\varepsilon}_{#1}
% ... following Federer's 1969 book 2.5.19
\newcommand{\Dirac}[1]{ \boldsymbol{\delta}_{#1} }

% Smooth compactly supported vectorfields
\newcommand{\VF}{\mathscr{X}}

% A symbol for the Dirac measure
\newcommand{\ellipticity}[1]{\operatorname{ellipticity}(#1)}

% affine approximate
% \DeclareMathOperator{\Aff}{A}       

% affine space generated by a set
\DeclareMathOperator{\aff}{af{}f}

% The trace
\DeclareMathOperator{\trace}{trace}

\DeclareMathOperator{\Tr}{Tr}

% The greatest commands divisor
% \DeclareMathOperator{\gcd}{gcd}

% The neighbours of a cube
% \DeclareMathOperator{\nbs}{nbs}

% Space of homomorphisms between vector spaces
\DeclareMathOperator{\Hom}{Hom}

% The boundary of a set
\DeclareMathOperator{\Bdry}{Bdry}

% something approximate
\DeclareMathOperator{\ap}{ap}

\DeclareMathOperator*{\aplim}{\mathrm{ap}\, \lim}   % approximate limit
\DeclareMathOperator*{\aplimsup}{\mathrm{ap}\, \limsup}   % approximate limit

% linear span
\DeclareMathOperator{\lin}{span}

% class of continuity
\newcommand{\cnt}[1]{\mathscr{C}^{#1}}

% skeleton of a simplicial complex
\newcommand{\skel}[2]{\mathrm{sk}_{#1}(#2)}

% the set of orthogonal projections "O^*(n,m)" from R^n onto R^m with bold "O",
% produced by '\orthproj{n}{m}'
\newcommand{\orthproj}[2]{\mathbf{O}^\ast({#1},{#2})}
\newcommand{\orthinj}[2]{\mathbf{O}({#1},{#2})}

\newcommand{\powerset}[1]{\mathbf{2}^{#1}}

% the orthogonal group O(n)
\newcommand{\orthgroup}[1]{\mathbf{O}({#1})}

% the H\"older semi norm of f with exponent \alpha is "h_{\alpha} (f)" with
% bold \alpha, '\hoelder{\alpha}{f}'
\newcommand{\hoelder}[2]{\mathbf{h}_{#1}(#2)}

% the evaluation map  \ev{X} : X --> X^{**}
\newcommand{\ev}[1]{\mathbf{ev}_{#1}}

% An unnamed parameter/placeholder
\newcommand{\ph}{\,\cdot\,}

% Interior of a set
\DeclareMathOperator{\Int}{Int}

% Relative interior of a set
\DeclareMathOperator{\rInt}{relint}

% Barrier cone
\DeclareMathOperator{\bc}{bc}

% Effective domain of a convex function
\DeclareMathOperator{\dom}{dom}

% Relative boundary of a set
\DeclareMathOperator{\rBdry}{relbd}

% domain
\DeclareMathOperator{\reach}{reach}

% domain
\DeclareMathOperator{\dmn}{dmn}

% gradient
\DeclareMathOperator{\grad}{grad}
% \DeclareMathOperator{\grad}{\nabla}

% subgradient of a convex function
% \newcommand{\subgrad}{\text{{\large $\nabla$}}}
% \newcommand{\subgrad}{\nabla}
\DeclareMathOperator{\subgrad}{subgr}

% support
\DeclareMathOperator{\spt}{spt}

% space generated be a k-vector
\DeclareMathOperator{\Sp}{space}

% convex hull
\DeclareMathOperator{\conv}{conv}

% uniform convex hull
\DeclareMathOperator{\uconv}{uc}

% cone over a set (as in the book of Rockafellar)
\DeclareMathOperator{\ray}{ray}

% convex cone over a set (as in the book of Rockafellar)
\DeclareMathOperator{\cone}{cone}

% tangent space
\DeclareMathOperator{\Tan}{Tan}

% normal space
\DeclareMathOperator{\Nor}{Nor}

% Lipschitz constant
\DeclareMathOperator{\Lip}{Lip}

% produces "card", cardinality of set
\DeclareMathOperator{\card}{card}

% diameter
\DeclareMathOperator{\diam}{diam}

% sign of
\DeclareMathOperator{\sgn}{sgn}

% % maximum
% \DeclareMathOperator{\max}{max}

% Set A without B, 'A \without B' produces "A \sim B"
\newcommand{\without}{\!\sim\!}
% \DeclareMathOperator{\without}{\sim}

\newcommand{\QEC}{\textrm{QEC}}
\newcommand{\EC}{\textrm{EC}}
\newcommand{\SAC}{\textrm{SAC}}
\newcommand{\SEC}{\textrm{SEC}}
\newcommand{\CI}{\textrm{CI}}
\newcommand{\USAC}{\textrm{USAC}}
\newcommand{\AC}{\textrm{AC}}
\newcommand{\AlmE}{\textrm{AE}}
\newcommand{\AlmUE}{\textrm{UAE}}

% classification { x : properties }, : is '\with'
\DeclareMathOperator{\with}{:}

% Subspace of End(R^n) with image in a fixed space
\newcommand{\imsp}{\mathcal{I}}

% image of a map
\DeclareMathOperator{\im}{im}

% boundary of a simplex
\DeclareMathOperator{\sBdry}{b}

% interior of a simplex
\DeclareMathOperator{\sInt}{i}

% barycentre of a simplex
\DeclareMathOperator{\baryc}{c}

% AUE
\DeclareMathOperator{\FAUE}{AUE}
% AE
\DeclareMathOperator{\FAE}{AE}
% BC
\DeclareMathOperator{\FBC}{BC}
% wBC
\DeclareMathOperator{\FwBC}{wBC}
% AC
\DeclareMathOperator{\FAC}{AC}

% To mark parapraph which lack final verification
\newcommand{\verify}{\framebox{{\color{Red} \textsc{verify}}~}}
\newcommand{\todo}{\framebox{{\color{Red} \textsc{to do}}}}
\newcommand{\maciej}[1]{\begin{center}\framebox{\begin{minipage}{0.95\textwidth} {\color{Fuchsia}\textsc{@ML:} \textsf{#1}} \end{minipage}}\end{center}}
\newcommand{\comment}[1]{{\color{teal}$\boldsymbol{\gg}$ \textsf{#1} $\boldsymbol{\ll}$}}

\newcommand{\cred}{\color{red}}
\newcommand{\cg}{\color{green}}
\newcommand{\cm}{\color{magenta}}


% Local Variables:
% coding: utf-8
% eval: (custom-set-variables  '(LaTeX-indent-level 4) '(LaTeX-item-indent -4) '(indent-tabs-mode nil))
% eval: (ispell-change-dictionary "british")
% mode: latex
% eval: (flyspell-mode)
% End:
